\section{Method}
\label{sec:method}

\subsection{ShortcutScore: Gradient-Structure Detection}

\noindent\textbf{Why Gradient-Structure Detection.}
Standard empirical risk minimization treats all training samples equally, optimizing the average token-level loss across the dataset. This is problematic because shortcut reasoning samples---those that achieve low loss by exploiting keyword correlations, answer memorization, or surface pattern matching---generate strong gradients that efficiently reduce training loss but do not support generalizable reasoning. Critically, such samples are often correctly labeled, making them invisible to conventional data quality filters or scalar loss-based criteria. We therefore detect shortcut-promoting samples from the geometry of their gradients rather than from labels or loss values, identifying two complementary geometric signatures: (1) \textbf{Non-Transfer Alignment.} per-sample gradients are poorly aligned with the validation gradient, signaling updates that reduce loss without improving held-out reasoning; (2) \textbf{Answer-Gradient Concentration.} gradient energy concentrates on final-answer tokens rather than intermediate reasoning tokens, signaling the model learns \emph{what to answer} rather than \emph{how to reason}. The detector is intentionally a \emph{rank} signal: it does not need to be a calibrated estimator of shortcut probability---only to order the samples well enough that the downstream surgery's per-sample conflict gate has the right working set to act on.

\noindent\textbf{Step 1: Non-Transfer Gradient Alignment.}
For each training sample $s = (x, y)$, we compute the per-sample gradient $\mathbf{g}_s = \nabla_\theta \ell(\theta; s)$ and the validation gradient
\begin{equation}
    \mathbf{g}_{\mathcal{V}} = \nabla_\theta \frac{1}{|\mathcal{V}|} \sum_{v \in \mathcal{V}} \ell(\theta; v).
\end{equation}
The non-transfer gradient alignment is then defined as:
\begin{equation}
    A(s) = \cos(\mathbf{g}_s,\, \mathbf{g}_{\mathcal{V}}) = \frac{\mathbf{g}_s^\top \mathbf{g}_{\mathcal{V}}}{\|\mathbf{g}_s\|\|\mathbf{g}_{\mathcal{V}}\|}.
\end{equation}
A low value of $A(s)$ indicates that the gradient of sample $s$ points in a direction inconsistent with improving validation reasoning performance---a hallmark of shortcut exploitation.

\noindent\textbf{Step 2: Answer-Gradient Concentration.}
Let $T_{\mathrm{ans}}$ and $T_{\mathrm{reason}}$ denote the index sets of answer tokens and reasoning tokens in the output sequence. The answer-gradient concentration is defined as
\begin{equation}
    R(s) = \frac{\sum_{t \in T_{\mathrm{ans}}} \|\nabla_\theta \ell_t\|}{\sum_{t \in T_{\mathrm{ans}} \cup T_{\mathrm{reason}}} \|\nabla_\theta \ell_t\|},
\end{equation}
where $\ell_t$ is the token-level loss at position $t$. A high value of $R(s)$ indicates that the model's parameter updates are driven predominantly by answer prediction rather than by learning intermediate reasoning steps. When a token-level answer/reasoning split is unavailable (answer-only supervision, $T_{\mathrm{reason}}{=}\emptyset$), we set $R(s){=}\tau_R$ so that $S(s)$ degenerates gracefully to an alignment-only score; Appendix~\ref{sec:appendix_graceful} gives the derivation.

\noindent\textbf{Step 3: ShortcutScore and Selection Set.}
The two signals combine into the \textbf{ShortcutScore}:
\begin{equation}
    S(s) = \alpha \cdot \max(0,\, \tau_A - A(s)) + \beta \cdot \max(0,\, R(s) - \tau_R),
\end{equation}
where $\tau_A$ and $\tau_R$ are thresholds defining low alignment and high concentration, and $\alpha, \beta$ are balancing hyperparameters (defaults in Appendix~\ref{sec:appendix_sensitivity}). Samples with high $S(s)$ are flagged as shortcut-promoting, forming the selection set
\begin{equation}
    \mathcal{S}^\star = \{\, s \in \mathcal{D} : S(s) > \tau_S^\star \,\},
\end{equation}
where $\tau_S^\star$ is set to the F1-optimal threshold on $S(s)$ over a held-out probe (in our experiments $\tau_S^\star \approx 0$, so the selection set is approximately the upper half of the score distribution; see \S\ref{sec:exp_detection}). Samples in $\mathcal{S}^\star$ enter the gradient-surgery step (\S\ref{sec:gradient_surgery}); samples outside $\mathcal{S}^\star$ contribute their gradient unchanged. The validation gradient $\mathbf{g}_{\mathcal{V}}$ is recomputed periodically to keep the alignment signal up to date without excessive overhead. Algorithmic details appear in Appendix~\ref{sec:appendix_algorithm}.


%-----------------------------------------------------------
\subsection{Conflict-Only Gradient Surgery}
\label{sec:gradient_surgery}

\noindent\textbf{Why Conflict-Only Surgery, Not Reweighting or Removal.}
Detection alone identifies which samples need intervention but does not modify a single update. Two natural correction strategies fail. (1) \textbf{Hard Removal.} Dropping the entire sample (e.g., data filtering or influence-based selection) discards the validation-aligned component the sample still carries. A shortcut-flagged sample is a \emph{mixed} object: even when its gradient direction is overall harmful, a non-trivial subspace of its gradient remains aligned with $\mathbf{g}_{\mathcal{V}}$, and that subspace would have helped training had it been preserved. (2) \textbf{Scalar Reweighting.} Multiplying the per-sample gradient by a continuous weight (e.g., $w(s) = \exp(-\lambda S(s))$) shrinks magnitude but leaves direction unchanged: a downweighted-but-still-pointing-the-wrong-way gradient is still pointing the wrong way. Empirically, scalar reweighting is decoupled from the detector's calibration and offers no gain over conflict-only surgery alone (see ablation in \S\ref{sec:exp_ablation}). The minimal correction that both (i) preserves the useful component of a shortcut sample and (ii) cancels its harmful direction is therefore \emph{conditional directional projection}: act on direction, not on magnitude, and act only when there is a genuine conflict to resolve. We instantiate this idea via a PCGrad-style~\cite{yu2020gradient} conflict-only projection, augmented with an optional answer-gradient suppression term where the second detection signal $R(s)$ indicates the gradient is dominated by answer tokens.

\noindent\textbf{Step 1: Conflict-Only Projection for Non-Transfer Alignment.}
For each $s \in \mathcal{S}^\star$, we remove the component of $\mathbf{g}_s$ that conflicts with $\mathbf{g}_{\mathcal{V}}$, following the PCGrad conflict-removal principle. Concretely, we project only when the two gradients genuinely disagree:
\begin{equation}
\label{eq:gradient_surgery}
    \mathbf{g}_s' =
    \begin{cases}
        \mathbf{g}_s - \gamma\,\dfrac{\mathbf{g}_s^\top \mathbf{g}_{\mathcal{V}}}{\|\mathbf{g}_{\mathcal{V}}\|^2}\,\mathbf{g}_{\mathcal{V}}, & \text{if } \mathbf{g}_s^\top \mathbf{g}_{\mathcal{V}} < 0, \\[4pt]
        \mathbf{g}_s, & \text{otherwise,}
    \end{cases}
\end{equation}
where $\gamma \in [0,1]$ is the projection strength ($\gamma{=}1$ fully removes the conflicting component). This conflict-gated rule is deliberately asymmetric: even a shortcut-flagged sample may carry a small positive alignment with $\mathbf{g}_{\mathcal{V}}$, and erasing that component would discard a validation-improving direction---the opposite of our goal. Projection fires only in the truly harmful case ($\mathbf{g}_s$ points \emph{against} validation), and the resulting $\mathbf{g}_s'$ then has non-negative inner product with $\mathbf{g}_{\mathcal{V}}$.

\noindent\textbf{Step 2: Suppression for Answer-Gradient Concentration.}
For $s \in \mathcal{S}^\star$ with $R(s) > \tau_R$, we further decompose $\mathbf{g}_s'$ into answer-token and reasoning-token components and suppress the answer-dominant part:
\begin{equation}
    \mathbf{g}_s'' = (1 - \rho)\,\mathbf{g}_s'^{\,\mathrm{ans}} + \mathbf{g}_s'^{\,\mathrm{reason}},
\end{equation}
where $\rho \in [0,1]$ is the suppression coefficient. Setting $\rho{=}0$ disables this term, recovering pure conflict-only projection; setting $\rho{=}1$ removes the answer-token gradient contribution entirely, forcing the model to update from intermediate reasoning signals only. Empirically, $\rho{>}0$ is beneficial on tasks whose shortcut signal manifests in answer tokens but harms autoregressive generation in large language models (see \S\ref{sec:exp_llm_transfer}); we therefore use $\rho{=}0$ as the default in the LLM regime.

\noindent\textbf{Combined Update Rule.}
At each training step, the per-sample gradient enters surgery only if $s \in \mathcal{S}^\star$. The batch gradient is the unweighted sum of corrected per-sample gradients, $\mathbf{g}_{\mathrm{batch}} = \sum_{s \in \mathrm{batch}, s \in \mathcal{S}^\star} \mathbf{g}_s'' \;+\; \sum_{s \in \mathrm{batch}, s \notin \mathcal{S}^\star} \mathbf{g}_s$, which then drives a standard optimizer step. The validation gradient $\mathbf{g}_{\mathcal{V}}$ is refreshed every $k$ steps (Section~\ref{sec:exp_llm_transfer}) so the conflict test and projection direction remain accurate. Pseudocode for the combined SART step is in Appendix~\ref{sec:appendix_algorithm}.

\paragraph{Scaling to pretrained LLMs.}
A naive implementation stores a dense per-sample gradient of dimension $|\theta|$, which is $\sim$28\,GB in FP32 for a 7B backbone. We keep SART tractable on pretrained LLMs by (i) restricting trainable capacity to LoRA adapters~\cite{hu2022lora} so each per-sample gradient collapses to $\sim$32\,MB without altering Eqs.~(5) or~(\ref{eq:gradient_surgery}); (ii) optionally replacing gradients by fixed-seed Johnson--Lindenstrauss sketches of dimension $d_{\mathrm{sk}}{=}1024$ during scoring, which preserves the cosine alignment $A(s)$ up to $\varepsilon$-error; and (iii) refreshing $\mathbf{g}_{\mathcal{V}}$ every $k$ steps rather than every step, since it changes slowly relative to $\theta$. These three mechanisms are orthogonal to the \textsc{ShortcutScore} and surgery formulations---any alternative inner-product-preserving proxy (Fisher-diagonal, last-layer gradients, influence approximations~\cite{koh2020understandingblackboxpredictionsinfluence}) can be slotted in without changing the method. Full derivations and ablations appear in Appendix~\ref{sec:scalable_gradients}.
